%==============================================================================
%  MATH 231 -- Linear Algebra I (Queens College, Fall 2026)
%  COMPUTATIONAL UNIT, PART 5 -- Matrix Multiplication as an Algorithm
%
%  Section numbering runs CONTINUOUSLY across the parts of the computational
%  unit:
%
%      Part 1  MATH231_computational_1_operation_counting        Sections 1-4
%      Part 2  MATH231_computational_2_k_nearest_neighbours      Sections 5-10
%      Part 3  MATH231_computational_3_k_means_clustering        Sections 11-15
%      Part 4  MATH231_computational_4_gram_schmidt              Sections 16-20
%      Part 5  MATH231_computational_5_matrix_multiplication     Sections 21-23
%
%  Hence \setcounter{section}{20}: this part opens at Section 21.
%
%  STATUS: complete.  Sections 21-23.  A SHORT part by design -- a mini lesson,
%  not a full treatment.  The mathematics of the product belongs to the matrix
%  unit (Sec. 8), which defines the rule and proves its laws; this part only
%  asks what the rule COSTS.  Resist re-deriving the definition here.
%
%  AUDIENCE: distributed to students.  No instructor-facing apparatus.
%
%  THE COUNT.  c_ij needs p multiplications and p-1 additions = 2p-1 flops, and
%  there are mn entries, giving mn(2p-1) ~ 2mnp.  The square case is 2n^3, the
%  SAME leading term as Gram-Schmidt in Part 4 -- Sec. 23.3 points this out,
%  because two different algorithms landing on 2n^3 for two different structural
%  reasons is worth a student noticing.
%
%  Sec. 23.4 mentions that O(n^3) is not known to be optimal.  Keep it to a
%  Looking-ahead: naming Strassen is enough, and the exponent race is a rabbit
%  hole that would swamp a three-section part.
%
%  Compile with pdflatex (standard TeX Live packages only).
%==============================================================================
\documentclass[11pt]{article}

\usepackage[margin=1in]{geometry}
\usepackage{amsmath,amssymb,amsthm}
\usepackage[dvipsnames]{xcolor}
\usepackage{enumitem}
\usepackage{booktabs}
\usepackage{array}
\usepackage[hidelinks]{hyperref}
\usepackage{titlesec}
\usepackage{needspace}
\usepackage{tikz}
\usetikzlibrary{arrows.meta}

%------------------------------------------------------------------ style ----
\titleformat{\section}{\large\bfseries\sffamily}{\thesection.}{0.6em}{}
\titleformat{\subsection}{\normalsize\bfseries\sffamily}{\thesubsection.}{0.5em}{}
\titlespacing*{\section}{0pt}{14pt}{6pt}

\theoremstyle{plain}
\newtheorem{thm}{Theorem}[section]

\theoremstyle{definition}
\newtheorem{defn}{Definition}[section]
\newtheorem{ex}{Example}[section]
\theoremstyle{remark}
\newtheorem*{rmk}{Remark}
\newtheorem*{warn}{A common confusion}

\newcommand{\lookahead}{\par\smallskip\noindent
  \textbf{\sffamily\textcolor{RoyalBlue}{Looking ahead.}}\ \ignorespaces}
\newcommand{\lookback}{\par\smallskip\noindent
  \textbf{\sffamily\textcolor{ForestGreen}{From the matrix unit.}}\ \ignorespaces}

%--------------------------------------------------------------- notation ----
\newcommand{\R}{\mathbb{R}}
\newcommand{\vc}[1]{\mathbf{#1}}
\newcommand{\mt}[1]{\mathbf{#1}}
\newcommand{\norm}[1]{\lVert #1\rVert}
\newcommand{\T}{^{\mathsf{T}}}

%--------------------------------------------------------------- pseudocode --
\newcommand{\kw}[1]{\textsf{\bfseries #1}}
\newcommand{\ind}{\hspace*{1.3em}}
\newcommand{\indd}{\hspace*{2.6em}}
\newcommand{\inddd}{\hspace*{3.9em}}
\newcounter{algoline}
\newenvironment{algo}[2]{%
  \par\medskip\needspace{8\baselineskip}%
  \begin{center}\begin{minipage}{0.94\textwidth}%
  \hrule height 0.9pt \vspace{3pt}%
  \noindent{\sffamily\bfseries #1}\par\vspace{2pt}%
  \noindent{\small #2}\par\vspace{3pt}%
  \hrule height 0.4pt \vspace{5pt}%
  \begin{list}{\textbf{\arabic{algoline}:}}%
    {\usecounter{algoline}\setlength{\leftmargin}{2.3em}%
     \setlength{\labelwidth}{1.8em}\setlength{\labelsep}{0.5em}%
     \setlength{\itemsep}{2pt}\setlength{\topsep}{0pt}%
     \setlength{\parsep}{0pt}}%
}{%
  \end{list}\vspace{5pt}\hrule height 0.9pt%
  \end{minipage}\end{center}\medskip%
}

\title{\vspace{-1.2cm}\textbf{\sffamily MATH 231 -- Linear Algebra I}\\[4pt]
       {\large\sffamily The Computational Unit, Part 5 \textbullet\
        Matrix Multiplication as an Algorithm}\\[2pt]
       {\normalsize\sffamily Kiely Hall 326 \textbullet\
        Instructor: Kevin Bedoya}}
\author{}
\date{}

\begin{document}
\setcounter{section}{20}  % this part opens at Section 21; see the header note
\maketitle
\vspace{-1.6cm}

\noindent\textbf{About this document.} A short part. The matrix unit defined the
matrix product and established its laws; none of that is repeated here. The
question here is narrower and purely computational: taken as an
\emph{algorithm}, what does that definition actually instruct a machine to do,
and how much work is it?

\begin{center}
\renewcommand{\arraystretch}{1.2}
\begin{tabular}{@{}l l l@{}}
\toprule
\textbf{Part} & \textbf{Sections} & \textbf{Subject} \\
\midrule
1 & \S1--\S4    & operation counting, Big-O, floating point, error \\
2 & \S5--\S10   & $k$-nearest neighbours \\
3 & \S11--\S15  & $k$-means clustering \\
4 & \S16--\S20  & the Gram--Schmidt process \\
5 & \S21--\S23  & matrix multiplication as an algorithm \\
\bottomrule
\end{tabular}
\end{center}

\medskip
\noindent\textbf{Assumed.} From the matrix unit: what a matrix is and the
index convention $a_{ij}$ (\S3.1); the definition of the matrix product and its
dimension requirement (\S8.1--\S8.2). From this unit: operation counting and
FLOPs (\S1, \S4.6) and Big-O (\S2.7).

\medskip
\noindent\textbf{What you should be able to do afterwards.} State the condition
two matrices must satisfy before their product exists and give the shape of the
result; write the pseudocode for the product and say what each of its three
loops is responsible for; and count its operations and give its complexity.

%==============================================================================
\section{The operation, reviewed}
%==============================================================================

%------------------------------------------------------------------------------
\subsection{What is being computed}
%------------------------------------------------------------------------------
\lookback For $\mt{A}\in\R^{m\times p}$ and $\mt{B}\in\R^{p\times n}$, the
product $\mt{C}=\mt{A}\mt{B}$ is the $m\times n$ matrix with
\begin{equation}\label{eq:mm}
  c_{ij} \;=\; \sum_{k=1}^{p} a_{ik}\,b_{kj}
\end{equation}
(\S8.1 of the matrix unit). Entry $c_{ij}$ is the dot product of row $i$ of
$\mt{A}$ with column $j$ of $\mt{B}$.

Read \eqref{eq:mm} as an instruction rather than a description and it says: to
fill in one entry of the answer, walk along a row of $\mt{A}$ and down a column
of $\mt{B}$ in step, multiplying each pair and accumulating. Do that once for
every position in the output grid.

That is already the algorithm. Everything below is bookkeeping around it.

%------------------------------------------------------------------------------
\subsection{What must be true first}
%------------------------------------------------------------------------------
Before any arithmetic happens, the shapes must be checked.

\begin{center}
\fcolorbox{RoyalBlue!45}{RoyalBlue!7}{%
\begin{minipage}{0.9\textwidth}
\centering
\textbf{Precondition.} $\mt{A}\mt{B}$ exists only if the number of
\textbf{columns of $\mt{A}$} equals the number of \textbf{rows of $\mt{B}$}.
\[
  (m\times \underbrace{p)\ (p}_{\text{must agree}}\times n)
  \;\longrightarrow\; m\times n
\]
\end{minipage}}
\end{center}

\noindent The reason is visible in \eqref{eq:mm}: the sum pairs $a_{ik}$ with
$b_{kj}$ for $k=1,\dots,p$, so the row of $\mt{A}$ and the column of $\mt{B}$
must have the same number of entries or some term has no partner. The shared
value $p$ is what the sum runs over.

An implementation that skips this check does not get a wrong answer --- it reads
memory that does not belong to it, and the failure is arbitrary. Verifying the
shapes is the first line of the procedure for that reason, not out of
tidiness.

\begin{ex}
\[
  \begin{array}{lll}
    (4\times3)(3\times5) & \to\ 4\times5   & \text{defined}\\
    (3\times5)(4\times3) & \to\ \text{---} & \text{undefined: } 5\neq4\\
    (2\times2)(2\times2) & \to\ 2\times2   & \text{defined}\\
    (1\times n)(n\times1) & \to\ 1\times1  & \text{defined --- one dot product}
  \end{array}
\]
The second line is the first with the operands swapped, and it does not exist.
Order is part of the problem statement.
\end{ex}

%------------------------------------------------------------------------------
\subsection{The size of the job}
%------------------------------------------------------------------------------
Three numbers describe the whole computation, and it is worth separating what
each one controls:
\begin{center}
\renewcommand{\arraystretch}{1.3}
\begin{tabular}{@{}c l@{}}
\toprule
\textbf{Quantity} & \textbf{What it governs} \\
\midrule
$m$ & the number of rows in the answer \\
$n$ & the number of columns in the answer \\
$p$ & the length of each dot product --- the work \emph{inside} one entry \\
\bottomrule
\end{tabular}
\end{center}

\noindent So $mn$ counts \emph{how many} entries must be produced, and $p$
counts \emph{how expensive each one is}. Note that $p$ does not appear in the
output at all: it is summed away, and two matrices with a huge shared dimension
can produce a tiny answer. A $(2\times10^{6})(10^{6}\times2)$ product returns a
$2\times2$ matrix after a great deal of work.

%==============================================================================
\section{The algorithm}
%==============================================================================

%------------------------------------------------------------------------------
\subsection{Pseudocode}
%------------------------------------------------------------------------------
\begin{algo}{Matrix multiplication}{%
  \textbf{Input:} $\mt{A}\in\R^{m\times p}$, $\mt{B}\in\R^{p\times n}$.\quad
  \textbf{Output:} $\mt{C}=\mt{A}\mt{B}\in\R^{m\times n}$.}
\item \kw{if} $\operatorname{cols}(\mt{A}) \neq \operatorname{rows}(\mt{B})$:
        \kw{stop} --- the product is undefined
\item \kw{for} $i = 1,\dots,m$:
        \hfill\textit{\small pick a row of the answer}
\item \ind \kw{for} $j = 1,\dots,n$:
        \hfill\textit{\small pick a column of the answer}
\item \indd $s \;\gets\; 0$
        \hfill\textit{\small start the accumulator}
\item \indd \kw{for} $k = 1,\dots,p$:
        \hfill\textit{\small walk the row and the column together}
\item \inddd $s \;\gets\; s + a_{ik}\,b_{kj}$
\item \indd $c_{ij} \;\gets\; s$
        \hfill\textit{\small one entry finished}
\item \kw{return} $\mt{C}$
\end{algo}

%------------------------------------------------------------------------------
\subsection{Reading the three loops}
%------------------------------------------------------------------------------
The structure is three nested loops, and each has a distinct job:

\begin{itemize}[leftmargin=1.5em,itemsep=4pt]
  \item The \textbf{outer two} ($i$ and $j$) do no arithmetic whatsoever. They
        exist only to name a position in the output --- together they visit each
        of the $mn$ entries of $\mt{C}$ exactly once.
  \item The \textbf{inner one} ($k$) does all the work. For the position
        currently selected, it evaluates the dot product in \eqref{eq:mm}.
\end{itemize}

\noindent So the algorithm is not really three loops deep in any conceptual
sense. It is a double loop over the answer's entries, with a dot product
performed at each stop. Line 6 is the only line that multiplies anything, and
counting how many times it runs is the entire cost analysis of \S23.

\begin{rmk}[The accumulator]
Lines 4--7 are the standard shape of a summation in code: set a running total to
zero, add one term at a time, store the result when the loop ends. The variable
$s$ holds a partial sum that is meaningless until the loop finishes --- only
after $k$ has run its full range does $s$ equal $c_{ij}$.
\end{rmk}

\begin{warn}[Order of the loops]
The three loops may be nested in any order without changing the answer, because
addition is commutative and every $c_{ij}$ accumulates the same $p$ terms
regardless of the sequence they arrive in. It does \emph{not} follow that the
orderings perform equally well on a real machine: they differ in the pattern by
which they walk through memory, and that pattern can change the running time by
a large factor. The operation count in \S23 is identical for all of them, which
is a reminder that an operation count is not a running time.
\end{warn}

%==============================================================================
\section{What it costs}
%==============================================================================

%------------------------------------------------------------------------------
\subsection{One entry}
%------------------------------------------------------------------------------
Fix a position $(i,j)$ and count what lines 4--7 do. The inner loop runs $p$
times, and each pass performs one multiplication and one addition. The very
first addition is to zero and could be skipped, so:
\begin{center}
\renewcommand{\arraystretch}{1.3}
\begin{tabular}{@{}l l@{}}
\toprule
\textbf{Operation} & \textbf{Count} \\
\midrule
multiplications $a_{ik}b_{kj}$ & $p$ \\
additions into $s$             & $p-1$ \\
\midrule
\textbf{per entry}             & $2p-1$ FLOPs \\
\bottomrule
\end{tabular}
\end{center}

\noindent This is the dot-product count from \S20.1, unchanged --- one entry of
a matrix product is one dot product of length $p$.

%------------------------------------------------------------------------------
\subsection{The total}
%------------------------------------------------------------------------------
There are $mn$ entries and each costs $2p-1$, and no entry's work is shared with
any other's. So
\[
  \text{total}
  \;=\; mn\,(2p-1)
  \;=\; 2mnp - mn
  \;\approx\; 2mnp \text{ FLOPs} .
\]

\begin{thm}[Cost of matrix multiplication]
Multiplying $\mt{A}\in\R^{m\times p}$ by $\mt{B}\in\R^{p\times n}$ by the
algorithm of \S22.1 takes $\approx 2mnp$ FLOPs, so it runs in $O(mnp)$ time.
\end{thm}

\noindent Every one of the three dimensions appears exactly once, and each for a
reason already identified in \S21.3: $m$ and $n$ because that many entries must
be filled, $p$ because that is the length of each dot product.

%------------------------------------------------------------------------------
\subsection{The square case}
%------------------------------------------------------------------------------
When all three dimensions are equal --- two $n\times n$ matrices, the case that
comes up most --- set $m=n=p$:
\[
  2mnp \;=\; 2n^{3},
  \qquad\text{so the cost is } O(n^{3}) .
\]
The cubic is easy to see directly from the pseudocode: three nested loops, each
running $n$ times, with constant work at the innermost level.

\begin{rmk}
Doubling $n$ multiplies the work by roughly eight. Two $1000\times1000$
matrices --- a million entries each, which is not a large problem --- already
require about $2\times10^{9}$ floating-point operations.
\end{rmk}

\begin{rmk}[The same $2n^{3}$ as Gram--Schmidt]
Part 4 found that orthonormalising $n$ vectors in $\R^{n}$ also costs about
$2n^{3}$ FLOPs. The agreement is not a coincidence in any deep sense, but it is
not an accident either: both algorithms spend their time computing dot products
of length $n$, and both perform on the order of $n^{2}$ of them. The structures
differ --- matrix multiplication does exactly $n^{2}$ dot products, while
Gram--Schmidt does about $n^{2}/2$ and pays for the rest in subtractions --- yet
they land on the same leading term. Where an algorithm's cost comes from is
often more informative than the count itself.
\end{rmk}

%------------------------------------------------------------------------------
\subsection{Is cubic the best possible?}
%------------------------------------------------------------------------------
A natural assumption is that $n^{3}$ is forced --- the answer has $n^{2}$
entries and each looks like it needs $n$ operations, so what could be saved?

It is not forced.

\lookahead There are algorithms that multiply $n\times n$ matrices in
asymptotically fewer than $n^{3}$ operations. The first, due to Strassen, does
it in about $n^{2.81}$ by computing a $2\times2$ product with seven
multiplications instead of the obvious eight and applying that trick
recursively. Later methods have pushed the exponent lower still, though the
practical ones remain close to Strassen. The exponent that is actually optimal
is not known --- it cannot be below $2$, since the output alone has $n^{2}$
entries that must each be written, and closing the gap between $2$ and the best
known value is a genuine open problem.

For this course, $O(n^{3})$ is the operative figure: it is what the
straightforward algorithm costs, it is what the standard library routines
effectively deliver at the sizes we will meet, and it is the number to reach for
when estimating whether a computation is affordable.

\end{document}
